\documentclass[a4paper,12pt]{article}

\begin{document}
\thispagestyle{empty}
\begin{center}
\Large\bfseries
On the Number of Prime Numbers less than a Given Quantity.\\[6 pt]
(Ueber die Anzahl der Primzahlen unter einer gegebenen Gr\"{o}sse.)\\[12 pt]
Bernhard Riemann\\[12 pt]
[Monatsberichte der Berliner Akademie, November 1859.]\\[12 pt]
Translated by David R. Wilkins\\[24 pt]
\large\mdseries
Preliminary Version: December 1998
\end{center}

\newpage
\setcounter{page}{1}


\title{On the Number of Prime Numbers less than a Given Quantity.\\
(Ueber die Anzahl der Primzahlen unter einer gegebenen Gr\"{o}sse.)}
\author{Bernhard Riemann\\
Translated by David R. Wilkins\\
Preliminary Version: December 1998}
\date{[Monatsberichte der Berliner Akademie, November 1859.]\\
Translation \copyright\ D. R. Wilkins 1998.}

\maketitle

I believe that I can best convey my thanks for the honour which
the Academy has to some degree conferred on me, through my
admission as one of its correspondents, if I speedily make use of
the permission thereby received to communicate an investigation
into the accumulation of the prime numbers; a topic which perhaps
seems not wholly unworthy of such a communication, given the
interest which \emph{Gauss} and \emph{Dirichlet} have
themselves shown in it over a lengthy period.

For this investigation my point of departure is provided by
the observation of \emph{Euler} that the product
\[ \prod \frac{1}{\displaystyle 1 - \frac{1}{p^s}}
   = {\textstyle\sum} \frac{1}{n^s},\]
if one substitutes for $p$ all prime numbers, and for $n$
all whole numbers.  The function of the complex variable~$s$
which is represented by these two expressions, wherever
they converge, I denote by $\zeta(s)$.  Both expressions converge
only when the real part of $s$ is greater than~$1$;
at the same time an expression for the function
can easily be found which always remains valid.  On making use of
the equation
\[ \int\limits_0^\infty e^{-nx} x^{s-1} \,dx = \frac{\Pi(s - 1)}{n^s} \]
one first sees that
\[ \Pi(s - 1) \zeta(s) = \int\limits_0^\infty \frac{x^{s-1}\,dx}{e^x - 1}.\]
If one now considers the integral
\[ \int\limits \frac{(-x)^{s-1} \,dx}{e^x - 1} \]
from $+\infty$ to $+\infty$ taken in a positive sense around a
domain which includes the value~$0$ but no other point of 
discontinuity of the integrand in its interior, then this
is easily seen to be equal to
\[ (e^{-\pi si} - e^{\pi si})
      \int\limits_0^\infty \frac{x^{s-1}\,dx}{e^x - 1},\]
provided that, in the many-valued function
$(-x)^{s-1} = e^{(s-1)\log(-x)}$,
the logarithm of $-x$ is determined so as to be real when $x$
is negative.  Hence
\[ 2 \sin \pi s \, \Pi(s - 1) \zeta(s)
   = i \int\limits_\infty^\infty \frac{(-x)^{s-1} \,dx}{e^x - 1},\]
where the integral has the meaning just specified.

This equation now gives the value of the function $\zeta(s)$ for
all complex numbers~$s$ and shows that this function is
one-valued and finite for all finite values of $s$ with the
exception of $1$, and also that it is zero if $s$ is equal
to a negative even integer.

If the real part of $s$ is negative, then, instead of being taken
in a positive sense around the specified domain, this integral can
also be taken in a negative sense around that domain containing all
the remaining complex quantities, since the integral taken though values
of infinitely large modulus is then infinitely small.  However,
in the interior of this domain, the integrand has discontinuities only
where $x$ becomes equal to a whole multiple of $\pm 2\pi i$, and the
integral is thus equal to the sum of the integrals taken in a negative
sense around these values.  But the integral around the value $n2\pi i$
is $= (-n2\pi i)^{s-1}(-2\pi i)$, one obtains from this
\[ 2 \sin \pi s \, \Pi(s - 1) \zeta(s)
   = (2\pi)^s {\textstyle\sum} n^{s-1} ((-i)^{s-1} + i^{s-1}),\]
thus a relation between $\zeta(s)$ and $\zeta(1 - s)$, which,
through the use of known properties of the function $\Pi$, may be
expressed as follows:
\[ \Pi \left( \frac{s}{2} - 1 \right) \pi^{-\frac{s}{2}} \zeta(s) \]
remains unchanged when $s$ is replaced by $1 - s$.

This property of the function induced me to introduce,
in place of $\Pi(s - 1)$, the integral
$\displaystyle \Pi \left( \frac{s}{2} - 1 \right)$ into the
general term of the series $\displaystyle \sum \frac{1}{n^s}$,
whereby one obtains a very convenient expression for the function
$\zeta(s)$.  In fact
\[ \frac{1}{n^s} \Pi \left( \frac{s}{2} - 1 \right) \pi^{-\frac{s}{2}}
   = \int\limits_0^\infty e^{-nn\pi x} x^{\frac{s}{2} - 1} \,dx,\]
thus, if one sets
\[ \sum_1^\infty e^{-nn\pi x} = \psi(x) \]
then
\[ \Pi \left( \frac{s}{2} - 1 \right) \pi^{-\frac{s}{2}} \zeta(s)
   = \int\limits_0^\infty \psi(x) x^{\frac{s}{2} - 1} \,dx,\]
or since
\[ 2 \psi(x) + 1 = x^{-\frac{1}{2}} \left(
         2\psi \left( \frac{1}{x} \right) + 1 \right),
   \mbox{ (Jacobi, Fund. S. 184)}\]
\begin{eqnarray*}
\Pi \left( \frac{s}{2} - 1 \right) \pi^{-\frac{s}{2}} \zeta(s)
   &=& \int\limits_1^\infty \psi(x) x^{\frac{s}{2} - 1} \,dx
      + \int\limits_1^\infty \psi \left( \frac{1}{x} \right)
            x^{\frac{s-3}{2}} \,dx \\
   & &+ {\textstyle\frac{1}{2}} \int\limits_0^1 \left(
            x^{\frac{s-3}{2}} - x^{\frac{s}{2} - 1}
            \right) \,dx \\
   &=& \frac{1}{s(s-1)} + \int\limits_1^\infty \psi(x) \left(
            x^{\frac{s}{2} - 1} + x^{-\frac{1+s}{2}}
            \right) \,dx.
\end{eqnarray*}

I now set $s = \frac{1}{2} + ti$ and
\[ \Pi \left( \frac{s}{2} \right) (s - 1) \pi^{-\frac{s}{2}} \zeta(s)
   = \xi(t),\]
so that
\[ \xi(t) = {\textstyle\frac{1}{2}}
      - (tt + {\textstyle\frac{1}{4}}) \int\limits_1^\infty
         \psi(x) x^{-\frac{3}{4}}
         \cos ({\textstyle\frac{1}{2}}t \log x) \,dx \]
or, in addition,
\[ \xi(t)
   = 4 \int\limits_1^\infty \frac{d(x^{\frac{3}{2}} \psi'(x))}{dx}
          x^{-\frac{1}{4}} \cos ({\textstyle\frac{1}{2}}t \log x) \,dx.\]

This function is finite for all finite values of $t$, and allows
itself to be developed in powers of $tt$ as a very rapidly converging
series.  Since, for a value of $s$ whose real part is greater
than $1$,
$\log \zeta(s) = - \sum \log (1 - p^{-s})$
remains finite, and since the same holds for the logarithms of the
other factors of $\xi(t)$, it follows that the function $\xi(t)$
can only vanish if the imaginary part of $t$ lies between
$\frac{1}{2} i$ and $-\frac{1}{2} i$.  The number of roots
of $\xi(t) = 0$, whose real parts lie between $0$ and $T$
is approximately
\[ = \frac{T}{2\pi} \log \frac{T}{2\pi} - \frac{T}{2\pi};\]
because the integral $\int\limits d\log \xi(t)$, taken in a positive
sense around the region consisting of the values of $t$ whose imaginary
parts lie between $\frac{1}{2} i $ and $-\frac{1}{2}i$ and
whose real parts lie between $0$ and $T$, is (up to a fraction
of the order of magnitude of the quantity $\frac{1}{T}$) equal to
$\displaystyle \left( T \log \frac{T}{2\pi} - T \right) i$;
this integral however is equal to the number of roots of
$\xi(t) = 0$ lying within in this region, multiplied by $2 \pi i$.
One now finds indeed approximately this number of real roots
within these limits, and it is very probable that all roots are real.
Certainly one would wish for a stricter proof here; I have
meanwhile temporarily put aside the search for this
after some fleeting futile attempts, as it appears unnecessary
for the next objective of my investigation.

If one denotes by $\alpha$ all the roots of the equation
$\xi(\alpha) = 0$, one can express $\log \xi(t)$ as
\[ {\textstyle\sum} \log \left( 1 - \frac{tt}{\alpha \alpha} \right)
      + \log \xi(0);\]
for, since the density of the roots of the quantity~$t$ grows
with $t$ only as $\displaystyle \log \frac{t}{2\pi}$, it follows
that this expression converges and becomes for an infinite~$t$ only
infinite as $t \log t$; thus it differs from $\log \xi(t)$ by
a function of $tt$, that for a finite~$t$ remains continuous
and finite and, when divided by $tt$, becomes infinitely small
for infinite~$t$.  This difference is consequently a constant,
whose value can be determined through setting $t = 0$.

With the assistance of these methods, the number of prime numbers
that are smaller than $x$ can now be determined.

Let $F(x)$ be equal to this number when $x$ is not exactly
equal to a prime number; but let it be greater by $\frac{1}{2}$
when $x$ is a prime number, so that, for any $x$ at which there
is a jump in the value in $F(x)$,
\[ F(x) = \frac{F(x + 0) + F(x - 0)}{2}.\]

If in the identity
\[ \log \zeta(s) = - {\textstyle\sum} \log (1 - p^{-s})
   = {\textstyle\sum} p^{-s}
     + {\textstyle\frac{1}{2} \sum} p^{-2s}
     + {\textstyle\frac{1}{3} \sum} p^{-3s} + \cdots \]
one now replaces
\[ p^{-s}  \mbox{ by } s\int\limits_p^\infty x^{-s-1}\,ds,\quad
   p^{-2s} \mbox{ by } s\int\limits_{p^2}^\infty x^{-s-1}\,ds,\ldots,\]
one obtains
\[ \frac{\log \zeta(s)}{s}
   = \int\limits_1^\infty f(x) x^{-s-1} \,dx,\]
if one denotes
\[ F(x) + {\textstyle\frac{1}{2}} F(x^{\frac{1}{2}})
        + {\textstyle\frac{1}{3}} F(x^{\frac{1}{3}}) + \cdots \]
by $f(x)$.

This equation is valid for each complex value $a + bi$ of $s$ for
which $a > 1$.  If, though, the equation
\[ g(s) = \int\limits_0^\infty h(x) x^{-s} \,d \log x \]
holds within this range, then, by making use of
\emph{Fourier}'s theorem, one can express the function~$h$ in
terms of the function~$g$.  The equation decomposes, if $h(x)$
is real and
\[ g(a + bi) = g_1(b) + i g_2(b),\]
into the two following:
\[ g_1(b)   =    \int\limits_0^\infty h(x) x^{-a} \cos (b \log x) \,d \log x,\]
\[ i g_2(b) = -i \int\limits_0^\infty h(x) x^{-a} \sin (b \log x) \,d \log x.\]

If one multiplies both equations with
\[ (\cos (b \log y) + i \sin (b \log y)) \,db \]
and integrates them from $-\infty$ to $+\infty$, then one obtains
$\pi h(y) y^{-\alpha}$ on the right hand side in both, on account
of \emph{Fourier}'s theorems; thus, if one adds both equations
and multiplies them by $iy^\alpha$, one obtains
\[ 2\pi i h(y) = \int\limits_{a - \infty i}^{a + \infty i} g(s) y^s \,ds,\]
where the integration is carried out so that the real part
of $s$ remains constant.

For a value of $y$ at which there is a jump in the value of
$h(y)$, the integral takes on the mean of the values of the
function~$h$ on either side of the jump.  From the manner in
which the function $f$ was defined, we see that it has
the same property, and hence in full generality
\[ f(y) = \frac{1}{2\pi i} \int\limits_{a - \infty i}^{a + \infty i}
      \frac{\log \zeta(s)}{s} y^s \,ds.\]

One can substitute for $\log \zeta$ the expression
\[ \frac{s}{2} \log \pi - \log (s - 1)
   - \log \Pi \left( \frac{s}{2} \right)
   + {\textstyle\sum}^\alpha \log \left(
         1 + \frac{(s - \frac{1}{2})^2}{\alpha\alpha} \right)
   + \log \xi(0) \]
found earlier; however the integrals of the individual terms
of this expression do not converge, when extended to infinity,
for which reason it is appropriate to convert the previous equation
by means of integration by parts into
\[ f(x) = - \frac{1}{2\pi i} \frac{1}{\log x}
         \int\limits_{a - \infty i}^{a + \infty i} \, \,
         \frac{\displaystyle d \frac{\log \zeta(s)}{s}}{ds}
         x^s \,ds \]

Since
\[ - \log \Pi \left( \frac{s}{2} \right)
   = \lim \left( \sum_{n=1}^{n = m}
         \log \left( 1 + \frac{s}{2n} \right)
         - \frac{s}{2} \log m \right),\]
for $m = \infty$ and therefore
\[ - \frac{\displaystyle d \frac{1}{s}
            \log \Pi \left( \frac{s}{2} \right)}{ds}
   = \sum_1^\infty \frac{\displaystyle d \frac{1}{s}
            \log \left( 1 + \frac{s}{2n} \right)}{ds},\]
it then follows that all the terms of the expression for $f(x)$, with
the exception of
\[ \frac{1}{2\pi i} \frac{1}{\log x}
         \int\limits_{a - \infty i}^{a + \infty i}
         \frac{1}{ss} \log \xi(0) x^s \,ds
   = \log \xi(0), \]
take the form
\[ \pm \frac{1}{2\pi i} \frac{1}{\log x}
         \int\limits_{a - \infty i}^{a + \infty i} \, \,
         \frac{\displaystyle d \left( \frac{1}{s} \log \left(
               1 - \frac{s}{\beta} \right) \right)}{ds}
         x^s \,ds.\]

But now
\[ \frac{\displaystyle d \left( \frac{1}{s} \log \left(
               1 - \frac{s}{\beta} \right) \right)}{d\beta}
   = \frac{1}{(\beta - s)\beta},\]
and, if the real part of $s$ is larger than the real part of
$\beta$,
\[ - \frac{1}{2\pi i} \int\limits_{a - \infty i}^{a + \infty i}
         \frac{x^s \,ds}{(\beta - s)\beta}
   = \frac{x^\beta}{\beta} = \int\limits_\infty^x x^{\beta - 1} \,dx,\]
or
\[ = \int\limits_0^x x^{\beta - 1} \,dx,\]
depending on whether the real part of $\beta$ is negative or
positive.  One has as a result
\[ \frac{1}{2\pi i} \frac{1}{\log x}
         \int\limits_{a - \infty i}^{a + \infty i} \, \,
         \frac{\displaystyle d \left( \frac{1}{s} \log \left(
               1 - \frac{s}{\beta} \right) \right)}{ds}
         x^s \,ds \]
\begin{eqnarray*}
   &=& - \frac{1}{2\pi i} \int\limits_{a - \infty i}^{a + \infty i}
         \frac{1}{s} \log \left( 1 - \frac{s}{\beta} \right)
         x^s \,ds \\
   &=& \int\limits_\infty^x \frac{x^{\beta - 1}}{\log x} \,dx
     + \mbox{ const.}
\end{eqnarray*}
in the first, and
\[
   =   \int\limits_0^x \frac{x^{\beta - 1}}{\log x} \,dx
     + \mbox{ const.} \]
in the second case.

In the first case the constant of integration is determined if one
lets the real part of $\beta$ become infinitely negative; in the
second case the integral from $0$ to $x$ takes on values separated
by $2\pi i$, depending on whether the integration is taken through
complex values with positive or negative argument, and becomes
infinitely small, for the former path, when the coefficient of $i$
in the value of $\beta$ becomes infinitely positive, but for the
latter, when this coefficient becomes infinitely negative.
From this it is seen how on the left hand side
$\displaystyle \log \left( 1 - \frac{s}{\beta} \right)$
is to be determined in order that the constants of integration
disappear.

Through the insertion of these values in the expression for $f(x)$
one obtains
\[ f(x) = {Li}(x) - {\textstyle\sum^\alpha} \left(
            {Li} \left( x^{\frac{1}{2} + \alpha i} \right)
            + {Li} \left( x^{\frac{1}{2} - \alpha i} \right)
            \right)
         + \int\limits_x^\infty \frac{1}{x^2 - 1} \frac{dx}{x \log x}
         + \log \xi(0),\]
if in $\sum^\alpha$ one substitutes for $\alpha$ all positive roots
(or roots having a positive real part) of the equation $\xi(\alpha) = 0$,
ordered by their magnitude.  It may easily be shown,
by means of a more thorough discussion of the function~$\xi$, that
with this ordering of terms the value of the series
\[ {\textstyle\sum} \left(
            {Li} \left( x^{\frac{1}{2} + \alpha i} \right)
            + {Li} \left( x^{\frac{1}{2} - \alpha i} \right)
            \right) \log x \]
agrees with the limiting value to which
\[ \frac{1}{2\pi i} \int\limits_{a - bi}^{a + bi} \, \,
      \frac{\displaystyle d \frac{1}{s} {\textstyle\sum}
         \log \left( 1 + \frac{(s - \frac{1}{2})^2}{\alpha \alpha}
         \right)}{ds} x^s \,ds \]
converges as the quantity $b$ increases without bound; however
when re-ordered it can take on any arbitrary real value.

From $f(x)$ one obtains $F(x)$ by inversion of the relation
\[ f(x) = {\textstyle\sum} \frac{1}{n}
         F \left( x^{\frac{1}{n}} \right),\]
to obtain the equation
\[ F(x) = {\textstyle\sum} (-1)^\mu \frac{1}{m}
         f \left( x^{\frac{1}{m}} \right),\]
in which one substitutes for $m$ the series consisting
of those natural numbers that are not divisible by any square
other than~$1$, and in which $\mu$ denotes the number of prime factors
of $m$.

If one restricts $\sum^\alpha$ to a finite number of terms, then the
derivative of the expression for $f(x)$ or, up to a part
diminishing very rapidly with growing $x$,
\[ \frac{1}{\log x} - 2 {\textstyle\sum^\alpha}
         \frac{ \cos (\alpha \log x) x^{-\frac{1}{2}}}{\log x} \]
gives an approximating expression for the density of the
prime number $+$ half the density of the squares of the prime
numbers $+$ a third of the density of the cubes of the prime
numbers etc.\ at the magnitude~$x$.

The known approximating expression $F(x) = {Li}(x)$ is therefore
valid up to quantities of the order $x^{\frac{1}{2}}$ and gives
somewhat too large a value; because the non-periodic terms in the
expression for $F(x)$ are, apart from quantities that do
not grow infinite with $x$:
\[ {Li}(x) - {\textstyle\frac{1}{2}}{Li}(x^{\frac{1}{2}})
           - {\textstyle\frac{1}{3}}{Li}(x^{\frac{1}{3}})
           - {\textstyle\frac{1}{5}}{Li}(x^{\frac{1}{5}})
           + {\textstyle\frac{1}{6}}{Li}(x^{\frac{1}{6}})
           - {\textstyle\frac{1}{7}}{Li}(x^{\frac{1}{7}})
           + \cdots \]

Indeed, in the comparison of ${Li}(x)$ with the number of prime
numbers less than $x$, undertaken by \emph{Gauss} and
\emph{Goldschmidt} and carried through up to $x =$ three
million, this number has shown itself out to be, in the first
hundred thousand, always less than ${Li}(x)$; in fact the
difference grows, with many fluctuations, gradually with $x$.
But also the increase and decrease in the density of the primes
from place to place that is dependent on the periodic terms
has already excited attention, without however any law governing
this behaviour having been observed.  In any future count
it would be interesting to keep track of the influence
of the individual periodic terms in the expression for the
density of the prime numbers.  A more regular behaviour than
that of $F(x)$ would be exhibited by the function $f(x)$,
which already in the first hundred is seen very distinctly
to agree on average with ${Li}(x) + \log \xi(0)$.

\end{document}

